% Part: first-order-logic
% Chapter: models-of-arithmetic
% Section: models-of-pa

\documentclass[../../../include/open-logic-section]{subfiles}

\begin{document}

\olfileid{mod}{mar}{mpa}
\section{$\Th{PA}$ના નિદર્શો}

\begin{explain}
$\Th{TA}$નો દરેક અપ્રમાણભૂત નિદર્શ $\Th{PA}$નો પણ અપ્રમાણભૂત નિદર્શ
છે. આપણે જાણીએ છીએ કે $\Th{TA}$ના અને તેથી $\Th{PA}$ના અપ્રમાણભૂત
નિદર્શો અસ્તિત્વ ધરાવે છે. આપણે એ પણ જાણીએ છીએ કે આવા અપ્રમાણભૂત
નિદર્શોમાં અપ્રમાણભૂત ``સંખ્યાઓ'' હોય છે, એટલે કે વ્યાપકક્ષેત્રના એવા
!!{element}s જે બધી પ્રમાણભૂત ``સંખ્યાઓ''ની ``પાર'' હોય છે. પરંતુ
તેઓ કેવી રીતે ગોઠવાયેલા છે? કેટલા છે? આપણે જોયું છે કે નિર્બળ
સિદ્ધાંત~$\Th{Q}$ના નિદર્શમાં માત્ર એક અપ્રમાણભૂત સંખ્યા પણ હોઈ શકે
છે. પરંતુ આવી સરળ !!{structure}s, $\Th{PA}$ કે $\Th{TA}$ના નિદર્શો
નથી.

$\Th{PA}$ કે $\Th{TA}$ના નિદર્શોની સંરચના સમજવાની ચાવી એ જોવામાં છે
કે આ સિદ્ધાંતોમાં કઈ હકીકતો !!{derivable} છે. દાખલા તરીકે, પહેલેથી
$\Th{PA}$, $\lforall[x][\eq/[x][x']]$ અને
$\lforall[x][\lforall[y][\eq[(x+y)][(y+x)]]]$ સાબિત કરે છે; તેથી જે
સરળ સંરચનાઓમાં આ !!{sentence}s અસત્ય છે તે $\Th{PA}$ના નિદર્શો બની
શકતી નથી.

ધારો કે $\Struct{M}$ એ $\Th{PA}$નો નિદર્શ છે. પછી જો
$\Th{PA}\Proves !A$, તો $\Sat{M}{!A}$. ફરી
$\Assign{\Obj{0}}{M}$ માટે $\nszero$, $\Assign{\prime}{M}$ માટે
$\nssucc$, $\Assign{+}{M}$ માટે $\nsplus$,
$\Assign{\times}{M}$ માટે $\nstimes$ અને $\Assign{<}{M}$ માટે
$\nsless$ વાપરીએ. ત્યારે કોઈ પણ !!{sentence}~$!A$, $\nszero$,
$\nssucc$, $\nsplus$, $\nstimes$ અને $\nsless$ વિશે કોઈ શરત કહે છે,
અને $\Sat{M}{!A}$ હોય તો એ શરત સંતોષાવી જોઈએ. દાખલા તરીકે, જો
$\Sat{M}{!Q_1}$, એટલે કે
$\Sat{M}{\lforall[x][\lforall[y][(\eq[x'][y']\lif\eq[x][y])]]}$,
તો $\nssucc$ !!{injective} હોવું જોઈએ.
\end{explain}

\begin{prop}
$\Struct{M}$માં $\nsless$ રેખીય કડક ક્રમ છે, એટલે તે નીચેની શરતો
સંતોષે છે:
\begin{enumerate}
\item કોઈ પણ $x\in\Domain{M}$ માટે $x\nsless x$ નહિ.
\item જો $x\nsless y$ અને $y\nsless z$, તો $x\nsless z$.
\item કોઈ પણ $x\neq y$ માટે $x\nsless y$ અથવા $y\nsless x$.
\end{enumerate}
\end{prop}

\begin{proof}
$\Th{PA}$ નીચેના વાક્યો સાબિત કરે છે:
\begin{enumerate}
\item $\lforall[x][\lnot x < x]$
\item $\lforall[x][\lforall[y][\lforall[z][((x < y \land y < z) \lif x < z)]]]$
\item $\lforall[x][\lforall[y][((x < y \lor y < x) \lor \eq[x][y])]]$
\end{enumerate}
\sourcecorrection{OLMAR-010}{સ્થિર અંગ્રેજી મૂળની
  \texttt{models-of-pa.tex}, પંક્તિ~56માં રેખીયતાના ત્રીજા સૂત્રના
  અંતે એક વધારાનું જમણું કૌંસ અને એક વધારું જમણું ચોરસ કૌંસ છે. અહીં
  બે પરિમાણક અને સૂત્રના બહારના એક કૌંસને બરાબર બંધ કરતાં સંતુલિત
  સૂત્રને પુનઃસ્થાપિત કર્યું છે.}
\end{proof}

\begin{prop}
\ollabel{prop:M-discrete} $\nszero$, $\nsless$-ક્રમમાં
$\Domain{M}$નો લઘુતમ !!{element} છે. દરેક~$x$ માટે
$x\nsless x^\nssucc$, અને $x^\nssucc$ આ ગુણધર્મ ધરાવતો
$\nsless$-લઘુતમ !!{element} છે. દરેક $x\neq\nszero$ માટે એવો એકમાત્ર
$y$ છે કે $y^\nssucc=x$. (આપણે $y$ને $\Struct{M}$માં $x$નો
``પૂર્વગામી'' કહીએ છીએ અને તેને $^\nssucc x$થી દર્શાવીએ છીએ.)
\sourcecorrection{OLMAR-011}{સ્થિર અંગ્રેજી મૂળની
  \texttt{models-of-pa.tex}, પંક્તિઓ~64--66 દરેક~$x$ને પૂર્વગામી છે
  એવો દાવો કરે છે, પરંતુ એ જ સિદ્ધાંતની $!Q_2$ સ્વયંસિદ્ધિ કહે છે કે
  $\nszero$ કોઈ ઉત્તરગામી નથી. અહીં દાવાને $x\neq\nszero$ સુધી મર્યાદિત
  કર્યો છે; નીચે પંક્તિઓ~111--113માં ખંડ બનાવતી વખતે પણ સાચી જરૂરી શરત,
  એટલે દરેક અપ્રમાણભૂત ઘટકને પૂર્વગામી છે, સ્પષ્ટ કરી છે.}
\end{prop}

\begin{proof}
અભ્યાસપ્રશ્ન.
\end{proof}

\begin{prob}
$\Th{PA}$માં !!{derivable} (અને તેથી $\Struct{N}$માં સત્ય) એવાં
$\Lang{L_A}$નાં !!{sentence}s શોધો જે
\olref[mod][mar][mpa]{prop:M-discrete}માં $\nszero$, $\nssucc$ અને
$\nsless$ માટે આપેલા ગુણધર્મોની ખાતરી આપે.
\end{prob}

\begin{prop}
$\Struct{M}$ના બધા પ્રમાણભૂત !!{element}s, બધા અપ્રમાણભૂત
!!{element}s કરતાં ($\nsless$ મુજબ) નાના છે.
\end{prop}

\begin{proof}
$\Struct{M}$ના પ્રમાણભૂત !!{element}~$\Value{\num{n}}{M}$ માટે આપણે
ટૂંકમાં $n$ લખીશું. પહેલેથી $\Th{Q}$ દરેક $n\in\Nat$ માટે
$\lforall[x][(x<\num{n}'\lif(\eq[x][\num{0}]\lor
\eq[x][\num{1}]\lor\dots\lor\eq[x][\num{n}]))]$ સાબિત કરે છે.
$\nsless\nszero$ હોય એવા કોઈ !!{element}s નથી. તેથી જો $n$
પ્રમાણભૂત અને $x$ અપ્રમાણભૂત હોય, તો $x\nsless n$ હોઈ શકે નહિ.
વ્યાખ્યા મુજબ અપ્રમાણભૂત ઘટક કોઈ પણ $n\in\Nat$ માટે
$\Value{\num{n}}{M}$ હોતો નથી, તેથી $x\neq n$ પણ. $\nsless$ રેખીય
ક્રમ હોવાથી $n\nsless x$ હોવું જ પડે.
\end{proof}

\begin{prop}
$\Domain{M}$નો દરેક અપ્રમાણભૂત !!{element}~$x$ નીચેના ઉપગણનો ઘટક છે:
\[
\dots ^{\nssucc\nssucc\nssucc}x \nsless ^{\nssucc\nssucc}x \nsless
^{\nssucc}x \nsless x \nsless x^{\nssucc} \nsless x^{\nssucc\nssucc}
\nsless x^{\nssucc\nssucc\nssucc} \nsless \dots
\]
આ ઉપગણને \emph{$x$નો ખંડ} કહીએ છીએ અને $[x]$ લખીએ છીએ. તેને લઘુતમ
કે મહત્તમ !!{element} નથી. તેને એવા $y\in\Domain{M}$ના ગણ તરીકે
લક્ષણિત કરી શકાય છે કે કોઈ પ્રમાણભૂત~$n$ માટે
$x\nsplus n=y$ અથવા $y\nsplus n=x$.
\end{prop}

\begin{proof}
આવો ગણ~$[x]$ હંમેશાં અસ્તિત્વ ધરાવે છે, કારણ કે દરેક અપ્રમાણભૂત
!!{element}~$y$ને $\Domain{M}$માં એકમાત્ર ઉત્તરગામી $y^\nssucc$ અને
એકમાત્ર પૂર્વગામી $^\nssucc y$ હોય છે. ક્રમિક !!{element}s $y$,
$y^\nssucc$ માટે
$y\nsless y^\nssucc$, અને $y^\nssucc$ એવો $\Domain{M}$નો
$\nsless$-લઘુતમ !!{element} છે કે $y$ તેનાથી $\nsless$-નાનો છે.
હંમેશાં $^\nssucc y\nsless y$ અને $y\nsless y^\nssucc$ હોવાથી
$[x]$ને લઘુતમ કે મહત્તમ !!{element} નથી. જો $y\in[x]$, તો
$x\in[y]$, કારણ કે ત્યારે કાં તો $y^{\nssucc\dots\nssucc}=x$ અથવા
$x^{\nssucc\dots\nssucc}=y$. જો $n$ $\nssucc$ સંકેતો સાથે
$y^{\nssucc\dots\nssucc}=x$, તો $y\nsplus n=x$, અને ઊલટું પણ,
કારણ કે $n$ એ $\prime$ સંકેતોની સંખ્યા હોય ત્યારે
$\Th{PA}\Proves\lforall[x][\eq[x^{\prime\dots\prime}][(x+\num{n})]]$.
\end{proof}

\begin{prop}
જો $[x]\neq[y]$ અને $x\nsless y$, તો દરેક $u\in[x]$ અને દરેક
$v\in[y]$ માટે $u\nsless v$.
\end{prop}

\begin{proof}
નોંધો કે
$\Th{PA}\Proves\lforall[x][\lforall[y][(x<y\lif(x'<y\lor x'=y))]]$.
તેથી જો $u\nsless v$ અને $[u]\neq[v]$, તો દરેક~$n$ માટે
$u\nsplus n^\nssucc\nsless v$ પણ.

દરેક $u\in[x]$ એ $\nsless y$ છે: પરિકલ્પનાથી $x\nsless y$. જો
$u\nsless x$, તો પરંપરિતતાથી $u\nsless y$. અને જો $x\nsless u$ અને
$u\in[x]$, તો કોઈ~$n$ માટે $u=x\nsplus n^\nssucc$, તેથી હમણાં સાબિત
કરેલી હકીકતથી $u\nsless y$.

હવે ધારો કે $v\in[y]$ એ $\nsless y$ છે, એટલે કોઈ પ્રમાણભૂત~$m$ માટે
$v\nsplus m^\nssucc=y$. ત્યારે $v\nsless x$ હોઈ શકે નહિ, નહિ તો
$y=v\nsplus m^\nssucc\nsless x$. સ્પષ્ટપણે $x\neq v$ પણ, નહિ તો
$x\nsplus m^\nssucc=v\nsplus m^\nssucc=y$ અને $[x]=[y]$ થાય. તેથી
$x\nsless v$. પરંતુ પછી દરેક~$n$ માટે
$x\nsplus n^\nssucc\nsless v$ પણ. આથી જો $x\nsless u$ અને
$u\in[x]$, તો $u\nsless v$. જો $u\nsless x$, તો પરંપરિતતાથી
$u\nsless v$.

છેલ્લે, જો $y\nsless v$, તો આપણે બતાવ્યું તેમ $u\nsless y$ અને
$y\nsless v$, તેથી $u\nsless v$.
\end{proof}

\begin{cor}
જો $[x]\neq[y]$, તો $[x]\cap[y]=\emptyset$.
\end{cor}

\begin{proof}
ધારો કે $z\in[x]$ અને $x\nsless y$. ત્યારે દરેક $u\in[y]$ માટે
$z\nsless u$. જો $z\in[y]$ હોત, તો $z\nsless z$ થાત. $y\nsless x$
હોય ત્યારે પણ એવી જ દલીલ લાગુ પડે છે.
\end{proof}

\begin{explain}
આનો અર્થ એ છે કે ખંડોને પોતાને $\nsless$ સાથે સુસંગત રીતે ક્રમ આપી
શકાય: $[x]\nsless[y]$ ત્યારે અને માત્ર ત્યારે જ્યારે $x\nsless y$;
અથવા સમકક્ષ રીતે, દરેક $u\in[x]$ અને $v\in[y]$ માટે $u\nsless v$.
સ્પષ્ટ છે કે પ્રમાણભૂત ખંડ~$[0]$ લઘુતમ ખંડ છે. તે કોઈ અપ્રમાણભૂત
ખંડને છેદતો નથી અને કોઈ બે અપ્રમાણભૂત ખંડો પણ એકબીજાને છેદતા નથી.
ખાસ કરીને, વારંવાર ઉત્તરગામી કે પૂર્વગામી લઈને જુદા ખંડ સુધી
``પહોંચી'' શકાય નહિ.
\end{explain}

\begin{prop}
જો $x$ અને $y$ અપ્રમાણભૂત હોય, તો $x\nsless x\nsplus y$ અને
$x\nsplus y\notin[x]$.
\end{prop}

\begin{proof}
$y$ અપ્રમાણભૂત હોય, તો $y\neq\nszero$.
$\Th{PA}\Proves\lforall[x][(y\neq\Obj{0}\lif x<(x+y))]$. હવે ધારો
કે $x\nsplus y\in[x]$. $x\nsless x\nsplus y$ હોવાથી
$x\nsplus n^\nssucc=x\nsplus y$ થાત. પરંતુ સરવાળાના રદ્દીકરણ નિયમ અનુસાર
$\Th{PA}\Proves\lforall[x][\lforall[y][\lforall[z][(\eq[(x+y)][(x+z)]\lif y=z)]]]$.
તેથી કોઈ પ્રમાણભૂત~$n$ માટે $y=n^\nssucc$ થાત, જ્યારે
$y$ અપ્રમાણભૂત હોવાનું ધાર્યું છે.
\end{proof}

\begin{prop}
લઘુતમ અપ્રમાણભૂત ખંડ નથી.
\end{prop}

\begin{proof}
$\Th{PA}\Proves\lforall[x][\lexists[y][(\eq[(y+y)][x]\lor
\eq[(y+y)'][x])]]$, એટલે દરેક~$x$ને~$2$ વડે ભાગી શકાય છે (સંભવતઃ
શેષ~$1$ સાથે). જો $x$ અપ્રમાણભૂત હોય, તો $y$ પણ અપ્રમાણભૂત છે. અગાઉના
પ્રસ્તાવથી $y\nsless y\nsplus y$ અને $y\nsplus y\notin[y]$. પછી
$y\nsless(y\nsplus y)^\nssucc$ અને
$(y\nsplus y)^\nssucc\notin[y]$ પણ. પરંતુ
$x=y\nsplus y$ અથવા $x=(y\nsplus y)^\nssucc$, તેથી $y\nsless x$ અને
$y\notin[x]$.
\end{proof}

\begin{prop}
મહત્તમ ખંડ નથી.
\end{prop}

\begin{proof}
અભ્યાસપ્રશ્ન.
\end{proof}

\begin{prob}
$\Th{PA}$ના અપ્રમાણભૂત નિદર્શમાં મહત્તમ ખંડ નથી તે બતાવો.
\end{prob}

\begin{prop}
\ollabel{prop:blocks-dense}
ખંડોનો ક્રમ સઘન છે. એટલે જો $x\nsless y$ અને $[x]\neq[y]$, તો બંનેથી
જુદો એવો તેમની વચ્ચેનો કોઈ ખંડ~$[z]$ છે.
\end{prop}

\begin{proof}
ધારો કે $x\nsless y$. અગાઉની જેમ $x\nsplus y$ને બે વડે ભાગી શકાય છે
(સંભવતઃ શેષ સાથે): એવો $z\in\Domain{M}$ છે કે કાં તો
$x\nsplus y=z\nsplus z$ અથવા
$x\nsplus y=(z\nsplus z)^\nssucc$. ઘટક~$z$, $x$ અને~$y$ની
``સરેરાશ'' છે, અને $x\nsless z$ તથા $z\nsless y$.
\sourcecorrection{OLMAR-012}{સ્થિર અંગ્રેજી મૂળની
  \texttt{models-of-pa.tex}, પંક્તિઓ~225--226માં આખા વિભાગ માટે
  જાહેર કરેલા નિદર્શ-સરવાળા $\nsplus$ને બદલે અજાહેર $\oplus$ અચાનક
  વપરાય છે. અહીં આ જ સાબિતીની આગલી પંક્તિ અને અગાઉની બધી સંજ્ઞા મુજબ
  ચારેય જગ્યાએ $\nsplus$ પુનઃસ્થાપિત કર્યું છે.}
\end{proof}

\begin{prob}
\olref[mod][mar][mpa]{prop:blocks-dense}ની વિગતવાર સાબિતી લખો.
$z$ના અસ્તિત્વની ખાતરી આપવા $\Th{PA}$એ કયું !!{sentence} !!{derive}
કરવું જોઈએ? $x\nsless z$ અને $z\nsless y$ શા માટે, અને
$[x]\neq[z]$ તથા $[z]\neq[y]$ શા માટે?
\end{prob}

\begin{explain}
જો $\Struct{M}$ !!{enumerable} હોય, તો તેના અપ્રમાણભૂત ખંડો સંમેય
સંખ્યાઓની જેમ ક્રમબદ્ધ હોય છે: તેઓ અંતબિંદુવિહીન !!a{denumerable}
સઘન રેખીય ક્રમ બનાવે છે. બતાવી શકાય છે કે આવા કોઈ પણ બે
!!{denumerable} ક્રમો એકરૂપ છે. તેથી સાચા અંકગણિતના કોઈ પણ બે
!!{enumerable} અપ્રમાણભૂત નિદર્શો $\Struct{M}_1$ અને
$\Struct{M_2}$ના માત્ર $<$ અને~$=$ ધરાવતી ભાષા સુધીના અવશેષો એકરૂપ
છે. ખરેખર, એકરૂપતા~$h$ આ પ્રમાણે વ્યાખ્યાયિત કરી શકાય: $\Struct{M_1}$
અને $\Struct{M_2}$ના પ્રમાણભૂત ભાગો પ્રમાણભૂત નિદર્શ~$\Struct{N}$ને,
અને તેથી એકબીજાને, એકરૂપ છે. અપ્રમાણભૂત ભાગ બનાવતા ખંડો પોતે સંમેય
સંખ્યાઓની જેમ ક્રમબદ્ધ છે અને તેથી એકરૂપ છે; ખંડોની એકરૂપતાને દરેક
ખંડની \emph{અંદર}ની એકરૂપતા સુધી વિસ્તારી શકાય—દરેક જોડીના ખંડોમાં
મનસ્વી ઘટકોને સામસામે ગોઠવો અને પછી $\Struct{M_1}$માં $x$ના
ઉત્તરગામીનું પ્રતિબિંબ, $\Struct{M_2}$માં $x$ના પ્રતિબિંબના ઉત્તરગામી
તરીકે લો. એમાંથી $\mathfrak{M}_1$ અને $\mathfrak{M}_2$ અંકગણિતની
પૂર્ણ ભાષામાં એકરૂપ છે એવું \emph{ફલિત થતું નથી} (વાસ્તવમાં,
એકરૂપતા હંમેશાં કોઈ !!a{language} સાપેક્ષ હોય છે), કારણ કે
$\Domain{M_1}$ અને $\Domain{M_2}$ પર સરવાળો અને ગુણાકાર વ્યાખ્યાયિત
કરવાની પરસ્પર બિનએકરૂપ રીતો છે. (વૉટના એક પ્રસિદ્ધ પ્રમેયમાંથી પણ આ
ફલિત થાય છે: કોઈ પૂર્ણ સિદ્ધાંતના ગણનીય નિદર્શોની સંખ્યા~2 હોઈ શકે
નહિ.)
\sourcecorrection{OLMAR-013}{સ્થિર અંગ્રેજી મૂળની
  \texttt{models-of-pa.tex}, પંક્તિઓ~238--242માં કોઈ ગણનીયતા-પરિકલ્પના
  વિના બધા અપ્રમાણભૂત ખંડોનો ક્રમ !!{denumerable} કહેવાયો છે. મનસ્વી
  નિદર્શ અગણનીય હોઈ શકે છે. અહીં આ દાવા માટે જરૂરી
  $\Struct{M}$ !!{enumerable} હોવાની શરત સ્પષ્ટ ઉમેરી છે; ત્યાર પછીનો
  બે !!{enumerable} નિદર્શો વિશેનો નિષ્કર્ષ યથાવત્ રહે છે.}
\end{explain}
\end{document}
